% Append-only section for the AION technical record.
% Requires amsmath; no additional packages are needed.
\section{Exact Shared-Prefix KV Reuse: Measured Boundary}
\label{sec:exact-prefix-kv-boundary}

\subsection{Question and hypothesis}

The experiment tested whether requests sharing an identical governed system,
Business Map, and response-contract prefix could reuse the corresponding neural
key--value (KV) state. The intended optimization was to calculate the common
prefix once on Apple Metal and process only each request's unique suffix. This
is genuine causal-state reuse, not semantic matching and not reuse based merely
on similar expert routes.

\subsection{Integrity contract}

The cache address cryptographically bound the exact token prefix, all local
model artifacts, FP16 datatype, and MPS backend. Reuse was forbidden unless all
fields matched. The model and evidence remained on the mounted SD volume.
Neither prompt text nor logit vectors were persisted. Promotion required both
bit-identical generated token identifiers and bit-identical logits at every
generation step, as well as at least a ten-percent total-time improvement after
charging the candidate for its one-time shared-prefix build.

\subsection{Method}

Four synthetic business requests contained 181--193 input tokens. Their exact
longest common token prefix contained 166 tokens, leaving 15--27 unique tokens
per request. The control performed each complete prefill in one Metal call. The
candidate calculated the 166-token prefix once, created independently growing
cache views without copying their underlying Metal tensors, preserved absolute
position and cache indices explicitly, and then processed each unique suffix.
Each arm greedily generated eight tokens per request.

An earlier engineering attempt tried to deep-copy non-leaf Metal cache tensors
and was rejected by PyTorch before evidence was emitted. It is not treated as
an experimental result. The corrected, position-bound run produced the
canonical evidence artifact
\texttt{exact\_prefix\_kv\_mps\_position\_bound\_20260908T0012Z.json}, whose
canonical report SHA-256 is
\texttt{e7ad2824fbe27bfe194ca78c05dd4cc98dc53c73cfe674de8c3f9e974852e150}.

\subsection{Results}

The control required 4.8451 seconds in total. The candidate required 4.6543
seconds including the shared-prefix build, an observed improvement of 3.94\%.
Control per-request time was 1.1072 seconds at p50 and 1.5882 seconds at
nearest-rank p95. Candidate suffix execution, excluding the separately charged
shared build, was 1.0604 seconds at p50 and 1.1095 seconds at p95.

The performance gate failed because 3.94\% is below the declared 10\% threshold.
The integrity gate also failed. Three of four generated token streams matched;
one differed. None of the four requests produced bit-identical logits at every
step. The largest step-logit difference observed after divergence was
$36.0342$.

\subsection{Conclusion and claim boundary}

The candidate is \texttt{NOT\_PROMOTED} and is stopped on the exact-output
track. Although causal KV state is mathematically reusable, splitting a
one-pass FP16 Metal prefill changes the floating-point kernel shape and therefore
does not preserve the control's bit pattern on the tested software stack. The
experiment does not show that KV caching is useless: it shows that this form
cannot be advertised as an exact-output optimization relative to one-pass
prefill, and its measured speed gain was too small in this workload.

The strategy may be reconsidered only as an explicitly quality-gated service or
with a backend that proves bit-identical chunked and one-pass prefill. Exact
work should instead prioritize avoiding neural calls entirely where outputs are
already verified: governed replay, narrowly verified AtomSheets, and finite
Semantic Gateway transformations. Those paths remove repeated calculations
without changing neural arithmetic.

\section{Native INT8 Glyph-Block Candidate}
\label{sec:native-int8-glyph-block}

After stopping exact prefix reuse, the model-level search moved to physical
weight traffic. The installed PyTorch/Metal stack exposes a native weight-only
INT8 matrix operation. This permits calculation directly from compact expert
weights instead of restoring each expert to a full FP16 tensor first.

A verified layer-16 Granite decode route containing eight experts was tested
over 200 order-balanced samples. Symmetric per-output-channel INT8 weights and
FP16 scales reduced the representation from 37,748,736 bytes to 18,915,328
bytes, a 49.89\% reduction. Native packed calculation with resident weights was
28.41\% faster at p50 than FP16. A naive fault implementation that transferred
each scale and matrix separately was 24.01\% slower and was rejected. Retaining
the small scales on Metal and joining both matrices of each expert into one
contiguous Glyph-addressed block made the complete placement-and-execution path
50.00\% faster at p50 than FP16.

This result is not exact. Maximum layer-output error was 0.0036621 and mean
absolute error was 0.00036268. It is therefore labelled
\texttt{ADVANCE\_TO\_FULL\_MODEL\_QUALITY\_GATE}, not promoted. Canonical SD-backed
evidence is recorded in
\texttt{int8pack-glyph-block-layer16-real-route-v2.json}; its canonical report
SHA-256 is
\texttt{48537cbcaba494e2e1d4555558d7bb585a1f53134a1926dc81e35e4536a45357}.

The measured ratio projects the complete 6.04 GB decimal Granite expert library
to approximately 3.03 GB decimal, or 2.82 GiB, including per-row scales. This
may allow all 1,280 experts to remain inside the existing 3 GiB expert budget,
removing repeated expert eviction after initialization. This projection is not
a full-model result. The required next experiment is a content-addressed
32-layer build followed by frozen output-quality, token-agreement, throughput,
time-to-first-token, peak-memory and cold/warm tests against both FP16 and the
promoted exact compressed streaming control.

\section{Full-Model Packed Library and Storage-First Boot}
\label{sec:int8-storage-first-boot}

The complete derived library contains 1,280 experts across all 32 Granite
layers. Its logical size is 3,026,452,480 bytes (2.8186 GiB), compared with
6,040,094,720 verified source bytes: a measured 49.894\% reduction. Every source
and derived object is hash-bound, and the library was atomically exposed only
after all layers and experts existed. The manifest canonical SHA-256 is
\texttt{62b7768bd2062d668a2c2aefbf765137829ac34390e8cb1421532c61dc3c3d5b}.

The first full-resident mechanism test installed native INT8 experts into all
32 layers. Across four frozen prompt families and eight generated tokens per
prompt, all 32 token identifiers matched FP16. Median time to first token
improved by 10.69\%. Median generation increased from 8.235 to 8.425 tokens per
second, only 2.30\%, below the declared 10\% speed gate. Logits were non-exact.
Direct full-resident speed is therefore \texttt{NOT\_PROMOTED}. Active Metal
allocation nevertheless fell from 6,597,587,456 to 3,591,071,232 bytes.

The storage-first implementation then constructed the model without the 64
combined FP16 expert tensors. It loaded 226 shared tensors and installed the
INT8 Glyph blocks directly. FP16 experts were never materialized. Ready-state
active Metal allocation was 3,584,241,664 bytes, 45.67\% below FP16. Boot time
fell from 91.2310 to 43.6290 seconds, a 52.18\% improvement. Median generation
was 8.215 tokens per second, 0.25\% below FP16, and all 32 short-run tokens again
matched. Every declared storage-first gate passed. The canonical evidence
SHA-256 is
\texttt{adb6e1e6deb9826c818bd4c1253bc298ba7b98d1360bc969fcc4f53c4df8a9e5}.

The status is \texttt{ADVANCE\_TO\_LONG\_QUALITY\_GATE}, not production
promotion. The result establishes a functioning fully resident model with
roughly half the expert storage, 45.67\% less active Metal memory and 52.18\%
faster boot while retaining short-run tokens and throughput. It does not prove
exact logits, 64--256-token quality, randomized prompt robustness or task
accuracy. Those are mandatory next-stage gates.

\section{Long-Generation Quality Boundary}
\label{sec:int8-long-quality-boundary}

The same four prompt families were extended to 64 generated tokens each. The
all-layer INT8 model produced finite logits, increased median throughput from
9.921 to 10.213 tokens per second (2.95\%), and reduced median time to first
token by 20.32\%. It matched 166 of 256 autoregressive token positions
(64.84\%), below the declared 75\% gate. It also missed the 10\% speed gate and
is \texttt{NOT\_PROMOTED}. The report canonical SHA-256 is
\texttt{a7c8bf2e0a4ba3d169561766103620803884fdb8460b1386b8e6424e75129b00}.

Two complete FP16 expert layers can be retained while keeping the expert
representation at 2.9940 GiB. Retaining the first and last layers improved
median throughput by 10.50\% and TTFT by 24.44\%, but token agreement was only
67.19\%. A real-activation sensitivity audit then substituted one packed layer
at a time across all four prompt families. Layers 29 and 1 were most sensitive,
followed by 17, 28 and 24. The audit canonical SHA-256 is
\texttt{190345280c86024d7f72ab693eb535ed9ac5924e4da4a0376381404ac7aeba93}.
Retaining layers 29 and 1 under the same budget produced only 64.06\% token
agreement and was also \texttt{NOT\_PROMOTED}.

Whole-layer mixed precision is stopped under the 3 GiB ceiling. The successful
short-run capacity and boot result remains valid, but long-form quality is not
promoted. The next candidate must improve the quantizer or retain sparse
high-value expert weights. It must use separate frozen calibration and
evaluation cohorts, teacher-forced diagnostics that distinguish one-token
divergence cascades from widespread probability damage, and semantic task
scoring in addition to raw token agreement.

A three-percent FP16 residual shelf was also tested on the real layer-16 route.
It kept the projected library within 2.9925 GiB, but reduced maximum layer error
only 1.32\% and slowed native packed calculation by 116.90\% because of the
additional indexing and FP16 correction kernels. It is
\texttt{NOT\_PROMOTED} and stopped. The canonical report SHA-256 is
\texttt{ca2604039afc04d5daf28611df89c8c96312d6c4a55e9cbf2223194e2ffd668a}.
This falsifies sparse input-column correction for the current quantizer: its
error is too distributed to repay the added kernels.

\subsection{Teacher-forced probability diagnostic}

To distinguish intrinsic quantization damage from autoregressive cascade, both
models were then evaluated on the same 256 FP16-generated token histories. The
packed model preserved the FP16 top-one choice at 247 positions (96.48\%) and
achieved 96.33\% mean top-five overlap. Mean KL divergence was 0.003455 nats,
and mean additional negative log likelihood on FP16-selected tokens was only
0.006691 nats. All four declared probability gates passed. The canonical report
SHA-256 is
\texttt{1c732ec78c27f73e769275c314caeaa08c637282fafab937110ad51f48ba3930}.

This result shows that raw positional agreement after free-running divergence
substantially overstates the underlying probability damage. It is labelled
\texttt{ADVANCE\_TO\_SEMANTIC\_TASK\_GATE}, not promoted. The logits are not
exact, and a frozen held-out evaluation of completed answer quality is required
before the packed model can be described as useful for long-form work.

\subsection{Frozen semantic microtask gates}

Two separately committed semantic cohorts were tested. A strict JSON cohort
was unsuitable because both FP16 and INT8 passed only two of eight tasks,
although seven parsed outputs were identical and INT8 was 6.26\% faster. A
twelve-case multiple-choice cohort removed JSON-format ambiguity. FP16 passed
only six cases, below its required ten-case competence floor. INT8 passed seven,
agreed with FP16 on eleven of twelve decisions (91.67\%), and ran 18.31\%
faster. Both gates are \texttt{NOT\_PROMOTED}; a candidate cannot inherit
absolute semantic credibility from a weak control. The multiple-choice report
canonical SHA-256 is
\texttt{4105f6fa4203073525639bf0b6cf7aca042f0f179bbb0582e2c3047551a00ba1}.

The observations add relative non-inferiority evidence but do not prove broad
capability. A larger frozen cohort must first establish FP16 competence using a
choice-likelihood method that is independent of response formatting. Because
the method is being refined after these results, a new unseen holdout remains
mandatory for promotion.

\subsection{Direct choice-likelihood diagnostic}

The free-generation parser was then removed as a possible confound. On the
already-observed twelve-case cohort, the evaluator performed one forward pass
per question and compared the next-token scores for A, B, C and D directly. It
combined the probability mass of the bare and space-prefixed single-token
spellings for each choice. This eliminated generation length, commentary and
answer parsing from the measurement.

FP16 nevertheless remained at six correct answers out of twelve. INT8 produced
eight correct answers, and the two models agreed on ten of twelve choices
(83.33\%). Aggregate candidate forward time was 3.3656 seconds versus 4.1288
seconds for FP16, 18.49\% lower, although speed was not a promotion criterion in
this post-hoc diagnostic. The result is
\texttt{DIAGNOSTIC\_DID\_NOT\_CLEAR\_BASELINE}. Its canonical report SHA-256 is
\texttt{7caf487aec34af497257f815f02e6e45b3bfba0a029cdd6353a278767117cd49}.

This falsifies response formatting as the principal cause of the control's weak
score. The cohort and evaluation method had already been observed, so this
result cannot promote the candidate in any event. An unseen holdout will not be
spent on the inadequate control design. The next semantic gate must first prove
FP16 competence on a larger independently frozen benchmark; only after that
control gate passes may the packed INT8 model be evaluated. The capacity, boot
and teacher-forced probability results remain valid, but absolute long-form
semantic quality remains unproven.

\subsection{Native Metal INT4 direct-path sweep}

The next mechanism gate used Apple's native Metal INT4 packed matrix operation
on the same real layer-16 eight-expert route as the successful INT8 test. Three
affine quantizers used groups of 32, 64 and 128 weights. Relative to FP16, their
logical expert representations were 68.75\%, 71.875\% and 73.4375\% smaller.
Their median resident expert calculations were respectively 34.83\%, 36.04\%
and 36.70\% faster. Thus four-bit packed execution is both physically smaller
and computationally functional on the existing M3 Pro Metal backend.

The quality gate did not pass. Group 32 was the most accurate INT4 condition,
with maximum layer-output error 0.0159302. The matched INT8 control measured
0.00299072, making group-32 error 5.33 times larger than INT8 and exceeding the
predeclared four-times ceiling. Groups 64 and 128 increased maximum error to
0.0227661 and 0.0281067. The result is
\texttt{STOP\_INT4\_DIRECT\_PATH}; canonical report SHA-256:
\texttt{1f6a3b60c40fe351e635140c1ce0a31c4ee4854b65c85c206be8184055955b07}.

Naive affine INT4 therefore does not advance to a full-model build. Its storage
and speed measurements establish a useful frontier, not a deployable model.
The gate will not be weakened and the direct group-size sweep will not be
repeated. A future four-bit candidate must use activation-aware quantization,
freeze separate calibration and evaluation activations, and improve upon the
recorded group-32 error/speed frontier before scaling.

\subsection{Activation-aware INT4 clipping boundary}

A second four-bit mechanism gate replaced random activation with captured real
Granite layer-16 activation. Commercial-risk and technical-explanation prompt
families formed the calibration partition. Business-planning and cache-diagnosis
families were withheld for evaluation. For every input and output expert matrix,
the calibration-weighted reconstruction objective selected among clipping
ratios 1.00, 0.98, 0.95, 0.90 and 0.85. Seventy-nine of eighty matrices selected
0.98 and one selected 0.95.

On the disjoint evaluation activations, normalized RMSE improved only from
0.12036 for naive group-32 INT4 to 0.11881, and mean absolute error improved by
1.62\%. Maximum error worsened from 0.222656 to 0.310547. The matched INT8
maximum error was 0.0449219, so the activation-aware INT4 maximum was 6.91 times
larger. Median one-token layer time was still 19.46\% faster than FP16 and
slightly faster than naive INT4, but both declared error gates failed. The
result is \texttt{STOP\_ACTIVATION\_AWARE\_INT4\_V1}; canonical report SHA-256:
\texttt{2cb2f415902b2f802f769c215ac8a00f3f420c517166b5ee3fc12f335b4d5e78}.

This falsifies simple activation-weighted clipping as the missing four-bit
mechanism. It did not materially address distributed quantization error and
increased the worst error. Ratio tuning will not be repeated. Any later INT4
candidate must change the equivalent weight/activation transformation itself,
charge runtime rescaling explicitly, and win a new narrow gate before a
full-model library is constructed.

\subsection{Equivalent channel-rescaled INT4 boundary}

The final four-bit mechanism candidate changed the weight/activation geometry
rather than clipping. Each input channel of a stored weight matrix was
multiplied by a bounded calibration-derived scale, while the corresponding live
activation channel was divided by the same scale. These operations cancel
exactly before quantization. Both runtime inverse scales and their storage were
charged. The resulting representation, including FP16 scale vectors, remained
68.66\% smaller than FP16. A balanced scaling exponent of 0.5 was selected for
73 of 80 matrices by the calibration objective.

On the disjoint real evaluation activations, normalized RMSE fell from 0.12036
to 0.11346 and mean absolute error improved by 5.64\%. Maximum error nevertheless
rose from 0.222656 to 0.246094, 5.48 times the matched INT8 error. After charging
the two inverse-scale operations, median one-token layer speed was only 6.93\%
better than FP16, compared with 17.22\% for naive INT4; the transformed path was
12.43\% slower than naive INT4. All four acceptance checks failed. The result is
\texttt{STOP\_EQUIVALENT\_SCALED\_INT4\_V1}; canonical report SHA-256:
\texttt{6f8a13860f37b9c7a749f8b44941aa9e8ae8016e8a3004547cb3025e7b30e7dd}.

The current four-bit branch is therefore closed. Direct affine INT4,
activation-weighted clipping and equivalent channel rescaling each failed the
declared quality/speed boundary. A full INT4 library will not be built from
these methods. Native INT8 remains the packed candidate that has passed the
mechanism, capacity, storage-first boot and teacher-forced probability gates.

\subsection{Sparse INT8 expert-dispatch breakthrough}

Runtime inspection identified a pure dispatch defect in the successful packed
INT8 path. Each decode token selected eight experts per layer, but the module
still invoked the input and output Metal matrix operations for all 40 experts.
Thirty-two zero-token experts caused 64 empty calls per layer, or 2,048
avoidable Metal calls across all 32 layers for each generated token. The
candidate removed only empty calls; selected experts, packed weights,
activations and their calculation order were unchanged.

Five hundred balanced trials used the real layer-16 eight-expert route. Metal
matrix calls fell from 80 to 16 per layer-token, an 80\% reduction. Median layer
time fell from 0.900542 ms to 0.548188 ms (39.13\%), while nearest-rank p95 fell
from 1.97037 ms to 1.55921 ms (20.87\%). Candidate and control packed-INT8
outputs were bit-identical and maximum error was zero. All mechanism gates
passed. The result is
\texttt{ADVANCE\_TO\_FULL\_MODEL\_SPARSE\_DISPATCH\_GATE}; canonical report
SHA-256:
\texttt{0a9244267122b65a627805bac04b51a7659416fbb87b543013f56ed240ebdf11}.

This is an exact-within-INT8 dispatch optimization, not an additional quality
trade. The layer result cannot by itself establish end-to-end generation speed,
because attention, routing, shared weights and generation orchestration remain
unchanged. Promotion requires a storage-first 64-token multi-family ABBA gate
with bit-identical INT8 token identifiers and step logits.

\subsection{Full-model sparse-dispatch promotion}

The first full-model form skipped both empty Metal matrix calls and their empty
concatenation entries. A storage-first four-family, 64-token ABBA run improved
median throughput by 13.93\%, from 10.254 to 11.682 tokens/s, and improved p95
generation time by 12.53\%. All token identifiers matched. Nevertheless, one of
twelve repeat comparisons produced non-identical logits with maximum error
1.17578125. The other eleven comparisons, including the second candidate run
for that family, were bit-identical. The implementation was
\texttt{NOT\_PROMOTED}; canonical report SHA-256:
\texttt{1facb2f8687be315f6cfb6ba1e4ba94ea2977ce959bce317d0ec62985f041d2b}.

The corrected implementation continued to omit every zero-token matrix
operation but restored the original 40-entry empty/non-empty concatenation
topology. Its repeated layer gate remained bit-identical, with p50 improved
31.43\% and p95 improved 15.67\%. The layer report canonical SHA-256 is
\texttt{f93eb681e4167d9f74c3ce9e37ac87bd286247e24472412d525f5f448f416015}.

The same full-model ABBA protocol was then repeated without prompt-specific
warm-up. Across 1,024 measured generated tokens, median control throughput was
10.144 tokens/s and sparse-dispatch throughput was 11.264 tokens/s, an 11.04\%
improvement. P95 generation time improved 7.88\%, from 6.7465 to 6.2150 seconds.
Every generated token identifier and every generation-step logit was
bit-identical across all twelve comparisons; maximum logit error was zero.
Storage-first construction again skipped all 64 FP16 expert tensors. All gates
passed, producing \texttt{PROMOTE\_INT8\_SPARSE\_DISPATCH}. Canonical report
SHA-256:
\texttt{73a2057a10c8f04f38a41f92bec6ad6a240b8033c958c3ee761bb3570cb43dee}.

This is a promoted exact-within-INT8 runtime optimization. It does not convert
the lossy INT8 representation into FP16 equivalence, and the evidence remains
bounded to four fixed prompt families and 64-token ABBA generation. The next
stage requires longer generation and a larger unseen prompt cohort under the
same bit-exact packed-INT8 requirement.

\subsection{Unseen 128-token replication boundary}

Eight new prompt families were frozen in isolated commit
\texttt{0853d69f} before either condition ran. The next storage-first ABBA gate
doubled response length to 128 tokens and measured 4,096 generated tokens across
32 runs. Sparse dispatch retained a 10.03\% median throughput improvement, from
10.2175 to 11.2427 tokens/s, and improved p95 generation time by 7.46\%.

Integrity did not hold universally. One of sixteen candidate repetitions first
changed token identifier at generation step 43 and reached maximum step-logit
error 37.7129; its other candidate repetition and closing baseline for the same
family were exact. Separately, one unchanged all-40 baseline repetition kept
every token identifier but produced maximum logit error 0.634766. The other 22
of 24 repeat comparisons were bit-identical. The result is
\texttt{NOT\_PROMOTED}; canonical report SHA-256:
\texttt{989ec69183409a6b3c93c7f5d07c2c5538172d2b35ae2ef88d20bd1f221ac287}.

The promoted four-family 64-token boundary remains valid. The larger cohort
does not establish exact 128-token behavior. Baseline self-drift shows that
Metal repeatability must be diagnosed separately, but candidate divergence also
occurred and the exact gate will not be weakened. Deterministic-mode repetition
of the two exposed families is the next diagnostic; even if successful, a new
full frozen replication is required before extending the claim.

\subsection{Deterministic-Metal diagnostic}

PyTorch deterministic algorithms were enabled in a post-hoc 128-token ABBA
diagnostic restricted to the two families that exposed drift. The runtime
accepted the setting, but every one of six repeat comparisons produced
non-identical tokens and logits, including both closing unchanged all-40
baseline runs. Maximum step-logit error reached 43.46875.

The setting also imposed substantial cost. Median control throughput was 7.130
tokens/s, compared with 10.217 in the ordinary eight-family run, while sparse
dispatch produced 7.706 rather than 11.243 tokens/s. Its relative median gain
fell to 8.08\%, below the ten-percent gate, although p95 generation time improved
14.22\%. The diagnostic is \texttt{NOT\_PROMOTED}; canonical report SHA-256:
\texttt{4e50a445d556b02c167f82c00ac62630625bf7d3a1c875d929f7725d671fcf04}.

Deterministic-algorithm mode is stopped on this MPS path. It neither made the
relevant kernels repeatable nor preserved the promoted speed gain. The exact
claim remains bounded to the successful 64-token sparse-dispatch cohort.
Longer quality-gated evidence must report baseline self-repeatability separately
and must not describe tolerance or token agreement as bit-exactness.

\subsection{Bounded persistent decode-arena boundary}

The sparse INT8 path still allocated two concatenation buffers per layer-token.
A candidate therefore preallocated input and output Metal arenas for eight
routed decode assignments and reused them; larger prefill batches retained the
existing sparse concatenation. The projected 32-layer arena occupied only
1,310,720 bytes (1.25 MiB), and candidate output was bit-identical.

Performance regressed. Across 500 balanced real-route trials, median layer time
increased from 0.602104 to 0.616250 ms, a 2.35\% slowdown. Nearest-rank p95
increased from 1.48804 to 1.63779 ms, a 10.06\% slowdown. The explicit slice
copy commands cost more than the allocation and concatenation they replaced.
The result is \texttt{STOP\_DECODE\_ARENA}; canonical report SHA-256:
\texttt{3fb648e8b208c00daef7eebb5833163b66faa4051fc724d5822ef949f0437b25}.

The persistent arena is not scaled or enabled. Default arena capacity is zero,
and sparse concatenation remains the promoted packed-INT8 dispatch path.

\subsection{Cached resident weight-view boundary}

The next candidate cached lightweight references to every expert's already
resident packed-weight and scale views. This avoided repeated attribute lookup,
slicing and reshaping without duplicating weight storage or changing any matrix
operation. On 1,000 balanced trials of the real layer-16 route, median dispatch
time improved by 11.31\%, from 0.607209 ms to 0.538562 ms, and nearest-rank p95
improved by 1.65\%. Layer outputs were bit-identical. The narrow gate advanced
to full-model testing; canonical report SHA-256:
\texttt{f29f2f17583ec86d10f7546cf58abb1051d6ed814d9f2ca9df133ecba16c441d}.

The four-family, 64-token storage-first ABBA test kept promoted sparse dispatch
active in both arms and changed only dynamic versus cached view construction.
Across eight observations per condition, median throughput decreased from
10.9462 to 10.8475 tokens/s, a 0.90\% regression. P95 generation time worsened
by 0.05\%. Every generated token identifier matched, but generation-step logits
were not universally bit-identical; maximum observed error was 0.46875. The
result is \texttt{NOT\_PROMOTED}; canonical report SHA-256:
\texttt{948136d201fd319a36439960410de325df70f201461fdb33d7b33966d8aaa384}.

The microbenchmark effect did not survive full-model execution and the exact
integrity gate also failed. Cached weight views are therefore stopped as a
production optimization. The promoted sparse dispatcher remains unchanged.

\subsection{Reusable empty-view boundary}

The exact sparse dispatcher retains all 40 concatenation entries and therefore
creates 64 zero-length tensors per layer-token, a projected 2,048 creations per
32-layer generated token. A bounded candidate reused one input and one output
zero-length Metal view. It added no weight storage and preserved the original
concatenation topology.

On 1,000 balanced trials of the real layer-16 route, median time improved by
6.81\%, from 0.609751 ms to 0.568229 ms, and output was bit-identical. The
nearest-rank p95 nevertheless increased by 3.82\%, from 1.50133 ms to 1.55867
ms. This exceeded the predeclared maximum 3\% tail-latency regression. The
result is \texttt{STOP\_CACHED\_EMPTY\_VIEWS}; canonical report SHA-256:
\texttt{4d8a5ab1d4ac140740bf8d590f7bb0cda0c42d4a03d73e64589880bd71a947eb}.

The candidate is not scaled to the full model or enabled by default. It shows
that reducing Python tensor-object churn can improve median dispatch time, but
this implementation did not deliver acceptable tail stability. The promoted
sparse dispatcher remains unchanged.

\subsection{Route-specialized compilation boundary}

A mechanism probe compiled one real packed-INT8 eight-expert graph on the local
MPS stack and produced bit-identical output. Granite's production router cannot
form one full dynamic graph because it converts data-dependent expert counts to
a Python list. A bounded alternative would cache a small number of fixed-route
compiled graphs per layer and retain eager execution as fallback.

Three hash-bound observations supplied 16,256 layer-route samples and 6,889
unique route sets when counted independently by layer. Aggregate exact-repeat
fraction was 57.62\%, but those repetitions were distributed across many route
patterns. The most common fixed graph covered only 3.05\% at the median layer;
the best eight graphs per layer covered only 12.99\%. This failed the
predeclared 50\% median-coverage requirement. The result is
\texttt{STOP\_ROUTE\_SPECIALIZED\_COMPILE}; canonical report SHA-256:
\texttt{0f037de61a3efb5e845294820ccaad11a42088fa567ec5092e839e4d050496d7}.

A fixed-route graph dictionary is not built. The useful surviving direction is
dynamic compilation: remove the router's Python \texttt{tolist()} boundary, or
build a fused Metal router/expert kernel that consumes expert indices directly
without specializing the graph to each route set.

\subsection{Fused eight-expert Metal matvec advance}

The native packed-INT8 Metal operation requires two-dimensional inputs and
cannot batch eight dynamically selected weight matrices. AION therefore
runtime-compiled a Metal SIMD matrix-vector kernel that accepts the eight
already-selected resident expert buffers and performs all eight one-token
products in one launch. Using it for both expert projections reduced expert
matrix launches from 16 to two per layer-token, an 87.5\% reduction, without
duplicating weights.

A 1,000-trial balanced gate exercised the complete input projection, SiLU/gate
and output projection on a real layer-16 eight-expert route. Median time improved
from 0.412355 ms to 0.329604 ms, a 20.07\% reduction. Nearest-rank p95 improved
from 1.35450 ms to 0.642416 ms, or 52.57\%. Maximum additional error relative
to packed INT8 was 0.0000305176 and RMSE was 0.000000307926, both below the
frozen 0.00075 maximum-error ceiling. Every gate passed, producing
\texttt{ADVANCE\_FUSED\_METAL\_MATVEC\_TO\_MODEL\_INTEGRATION}. Canonical
report SHA-256:
\texttt{18b1ce06a354cb4b1e5c480bf2bb5f84e34ceceafebac5177b2a1aa2ffb2238f}.

This result establishes the kernel mechanism only. Its altered reduction order
means it belongs to the explicitly separate quality-gated changed-kernel track,
not the bit-exact track. Full-model integration must be limited to one-token,
eight-unique-expert decode with the promoted packed-INT8 fallback retained, then
pass frozen probability, token, latency, memory and semantic gates.

\subsection{Fused-kernel full-model boundary}

The integration activated fused execution only for one-token routes containing
eight unique expert assignments and retained the promoted sparse packed-INT8
fallback everywhere else. Both conditions shared one storage-first resident
model; the candidate added 1.25 MiB of bounded output buffers.

Across the four-family, 64-token ABBA gate, median throughput improved from
10.9755 to 11.5191 tokens/s, a 4.95\% gain, and p95 generation time improved by
5.40\%. The candidate nevertheless missed its frozen 5\% throughput floor by
0.05 percentage points. Integrity also failed: only two of four families kept
identical generated token sequences. Both candidate repeats diverged for the
technical-explanation and business-planning families, whereas every closing
sparse-INT8 control repeat remained exact. Maximum accumulated step-logit error
was 38.5405. The result is \texttt{NOT\_PROMOTED}; canonical report SHA-256:
\texttt{fe7f1f63669671ab02509b7c9187bc0b4d7cf5a12975f1b91781f72abb91346b}.

Fused Metal matvec v1 is stopped at the full-model boundary and remains disabled
by default. The narrow kernel's small local error is amplified by recurrent
autoregressive execution. Its measured speed gain is preserved as evidence but
is not a promotion. A future kernel requires materially lower accumulated error
under a new frozen narrow gate before model-scale repetition.

\subsection{Calibration-selected fused-layer boundary}

Four frozen calibration families measured each of the 32 layers independently
for one decode step. The 16 layers with the lowest maximum second-step logit
disturbance were frozen as layers 0, 2--9, 12--14, 17, 22--23 and 29. The
cryptographically bound selection-plan SHA-256 is
\texttt{941f7224eef566f293937306302b3ad76a0280cb22a314f5798ae1a1ad07085c}.

Four disjoint frozen holdout families then underwent 64-token ABBA generation.
Median throughput improved from 11.0729 to 11.6042 tokens/s, a 4.80\% gain, but
p95 generation time regressed by 3.46\%. The logistics-planning family changed
tokens in both fused repeats. Document governance, product analysis and systems
explanation retained token identity but all exceeded the 0.05 step-logit ceiling.
Every closing sparse-INT8 control was exact. Maximum step-logit error reached
30.3359. The result is \texttt{NOT\_PROMOTED}; canonical report SHA-256:
\texttt{302f634116f56606e0fa6a9babfe8a863f84b3056fd33f60d8ee9ee97af44386}.

Sensitivity-selected fusion is stopped. It did not prevent recurrent numerical
amplification or tail-latency regression. Progressively smaller subsets must not
be tuned post hoc against this holdout. Any future fused candidate must improve
the arithmetic agreement of the kernel itself under a new frozen design.

\subsection{GPU-resident dynamic expert-bank breakthrough}

The next architecture removed Granite's per-layer GPU-to-Python route boundary.
Top-eight indices and gates remain on Metal and address one contiguous 40-expert
bank; two SIMD kernels perform both projections and reduce gated outputs on the
device. A 1,000-trial real layer-16 gate improved p50 from 1.87092 ms to 1.18408
ms (36.71\%) and p95 from 2.96121 ms to 2.21871 ms (25.07\%). Maximum extra
error was 0.0000610352, below the 0.00075 ceiling. The report SHA-256 is
\texttt{c093f623b2e5f6fe1b142f001a9775c1939352672b3b364269104bcf30f121b4}.

The SD-native conversion produced 32 independently hashed layer banks containing
3,026,452,480 logical bytes (2.8196 GiB), with zero duplicate runtime weight
bytes. Manifest canonical SHA-256:
\texttt{5ccd5b7d44467bd193e6cb8fc39b95df888a6f285a9d2e68690c685f24e8bb27}.

In four-family 64-token storage-first ABBA generation, the native-bank control
reached 11.2886 tokens/s and dynamic-bank execution reached 24.4679 tokens/s,
a 116.75\% improvement. P95 generation time improved 56.43\%, from 6.0314 to
2.62775 seconds. Both dynamic repeats retained identical token sequences in all
four families, and every closing control was bit-exact. Maximum step-logit
difference was 2.17969. The mechanism decision is
\texttt{ADVANCE\_DYNAMIC\_BANK\_TO\_QUALITY\_GATES}; report SHA-256:
\texttt{a167b4218c07d1a396093ab65d39ff8c111bad5c69b5af62b29e9372a6e5ece9}.

This is a genuine model-scale speed breakthrough and the first result near the
30-token/s target on the ordinary 18 GB Mac. It is not yet a quality promotion:
the changed kernel must pass direct FP16 teacher-forced probability and frozen
semantic validation.

\subsection{Dynamic-bank quality qualification and bounded promotion}

Incremental teacher forcing supplied FP16, bank-native INT8 and dynamic-bank
INT8 with identical histories across 256 frozen positions. Native INT8 retained
97.66\% FP16 top-one agreement, while dynamic banking retained 97.27\%, a
reduction of only 0.39 percentage points. Dynamic top-five overlap was 95.86\%,
mean KL divergence was 0.003687 nats and additional NLL on FP16-selected tokens
was 0.008386 nats. Every relative-quality gate passed, and dynamic execution
reduced teacher-forced time by 55.76\%. The decision was
\texttt{ADVANCE\_DYNAMIC\_BANK\_TO\_FROZEN\_SEMANTIC\_GATE}; report SHA-256:
\texttt{50bfc1e9928cfac76236f65f9ba3b432ca8dac4c52aea7e4cd169d2b078a0204}.

The previously frozen twelve-case semantic suite then produced identical answers
for native and dynamic INT8 in all 12 cases. Both scored 7/12 correct and every
answer was parseable. On these short tasks, dynamic throughput improved from
6.2627 to 9.4227 generated tokens/s, a 50.46\% increase. All gates passed and the
result is \texttt{PROMOTE\_DYNAMIC\_BANK\_WITHIN\_VALIDATED\_BOUNDARY}; canonical
report SHA-256:
\texttt{225b5fef106cf01f201adaef53df4b56b32b4ad262fcb4149b558083a7ed021a}.

The promotion is deliberately bounded. AION has demonstrated a 2.8196 GiB
SD-native expert bank running Granite at 24.47 tokens/s on four 64-token families,
with unchanged generated token sequences, strong FP16-relative probability
preservation and complete answer preservation on the frozen semantic suite. The
absolute 7/12 task score prevents any claim of universal model competence. This
is a promotion of the execution architecture and relative preservation. Unseen
128--256-token, repeatability and broader competent-model validation remain the
next expansion gates.

\subsection{Unseen long-generation replication and extended promotion}

The first long-generation replication reused an eight-family cohort that had
been frozen before dynamic-bank development. Across 32 ABBA runs and 4,096
generated tokens, dynamic banking sustained 24.8073 tokens/s at 128 tokens per
run, versus 11.2547 tokens/s for the native bank. This was a 120.42\% median
throughput improvement, while p95 generation time improved by 55.95\%.
Dynamic candidate repetitions were token-identical for all eight families;
native repetitions were token-identical for six. Canonical report SHA-256:
\texttt{6a9e929ffd2d527437ee1204a382a42655ebb554068452f48d3ef470b87a8933}.

A disjoint eight-domain holdout was then frozen in commit \texttt{7113f1ad},
before either condition executed. Its file SHA-256 is
\texttt{68bad763a3d70ba735779d4d6168114a11a4609183245d9c5c83b627221ac385}.
The 256-token ABBA run measured 8,192 generated tokens. Dynamic banking
sustained 23.9612 tokens/s versus 10.6228 tokens/s for the native bank, a
125.56\% improvement. P95 generation time improved by 59.05\%. The candidate
rate declined only 3.41\% when generation length doubled from 128 to 256.
Dynamic and native repetitions were each token-identical for seven of eight
families. Canonical report SHA-256:
\texttt{cfb5823bb2aa58c273769a8430296db3eb40ffd1a61b9a30b212de3773a916b9}.

Free-running native and dynamic output was not identical for seven of the eight
256-token families. This was not hidden or relabelled as exactness. The
changed-kernel track therefore underwent an extended teacher-forced comparison
on 2,048 identical positions from the new holdout. Native INT8 retained
98.1934\% FP16 top-one agreement and dynamic banking retained 98.2422\%.
Dynamic top-five overlap was 96.7969\%, mean KL divergence was 0.002552 nats,
and its FP16-token NLL delta was $-0.002641$ nats. Every predeclared relative
quality gate passed, while dynamic teacher time improved by 56.23\%. Canonical
report SHA-256:
\texttt{bd4e56118cec96351c3ee4db5d0a762590826668b9e5de57a072e1cad08601f4}.

One final hash-bound decision combined the 128-token evidence, the disjoint
256-token evidence, the 2,048-position teacher gate and the prior twelve-case
semantic gate. All seven acceptance conditions passed. The result is
\texttt{PROMOTE\_DYNAMIC\_BANK\_LONG\_GENERATION\_BOUNDARY}; canonical report
SHA-256:
\texttt{2e772e66d43a28c35cdf6c85b3027a08f29a5e0db1ab6fe39ceac2f3c4569620}.

The extended claim remains bounded to this Granite MoE model, the measured
18~GB M3 Pro, and the resident 2.8196~GiB SD-native expert bank. It establishes
sustained approximately 24-token/s execution through 256 generated tokens,
strong FP16-relative probability preservation, and repeatability no worse than
the native Metal control. It does not establish bit-exact logits, identical
free-running wording, arbitrary larger-model performance, or continuous
weight streaming from SD at the resident-bank rate.
